\documentclass[11pt]{article}

\usepackage[margin=1in]{geometry}
\usepackage{amsmath,amssymb}
\usepackage[T1]{fontenc}
\usepackage{lmodern}
\usepackage{hyperref}
\usepackage{array}
\usepackage{booktabs}

\title{Full Mathematical Note on the Symmetric and Asymmetric Clausius Workflow}
\author{}
\date{\today}

\begin{document}
\maketitle

\begin{abstract}
This document gives a fully mathematical explanation of the Clausius workflow
used in \texttt{nuclear\_matter\_workflows}. It explains the physics, the
mathematics, and the numerical logic of both:
\begin{itemize}
\item the zero-temperature asymmetric Clausius constraint analysis, where the
five interaction parameters
\[
a_n,\qquad a_{pn},\qquad b_n,\qquad b_{pn},\qquad c
\]
are determined;
\item the finite-temperature fixed-$y$ quantum critical-point solver, which
uses those five fitted parameters to compute
\[
T_c,\qquad n_c,\qquad P_c.
\]
\end{itemize}

The note is designed to be didactic. It does not skip algebraic steps, and it
uses simple physics analogies whenever useful.
\end{abstract}

\tableofcontents

\section{What Problem Is Being Solved?}

The repository is building an effective equation of state for nuclear matter.
The workflow is split into two layers.

\begin{enumerate}
\item First, at zero temperature, the code calibrates the interaction so that
it reproduces empirical nuclear-matter properties.
\item Second, at finite temperature, the code uses that calibrated interaction
to look for the liquid-gas critical point.
\end{enumerate}

So the logic is:
\[
\text{fit the EOS at } T=0
\quad \Longrightarrow \quad
\text{use that EOS at } T>0.
\]

The zero-temperature stage is where the five interaction parameters are found.
The finite-temperature stage does \emph{not} refit those parameters; it only
uses them.

\section{The Physical Ingredients of the Model}

The model has three basic physical ingredients.

\subsection{Quantum Kinetic Energy}

Nucleons are fermions. Because of the Pauli principle, even at zero
temperature they cannot all sit in the same momentum state. Instead, they fill
a Fermi sphere up to a Fermi momentum $k_F$.

This is the same basic physics as electrons in a metal. Even at $T=0$, the
system has pressure because the momentum states are filled.

\subsection{Repulsion Through Excluded Volume}

At short distances, nucleons cannot overlap arbitrarily. In an effective
description, this is represented by reducing the available volume.

An easy analogy is a box full of hard balls. If the balls occupy some volume,
the remaining free volume is smaller, so the same number of particles behaves
as if its density were larger.

\subsection{Attraction Through a Mean Field}

Nucleons also attract each other at intermediate range. In a mean-field model,
this appears as a negative interaction energy that scales roughly like the
number of interacting pairs, hence like $n^2$.

An easy analogy is a spring network in which every particle feels an average
pull from the others. The Clausius model modifies this pull by a denominator
$1+cn$, which makes the attraction soften as density grows.

\section{Step 1: The Zero-Temperature Ideal Fermi Gas}

\subsection{Density and Fermi Momentum}

At $T=0$, all one-particle momentum states are filled up to $k_F$. For a gas
with degeneracy $d$, the number density is
\begin{equation}
n_{\mathrm{id}}
=
\frac{d}{(2\pi)^3}\int_{|\mathbf{k}|\le k_F} d^3k
=
\frac{d}{2\pi^2}\int_0^{k_F} k^2\,dk
=
\frac{d\,k_F^3}{6\pi^2}.
\end{equation}

Solving for $k_F$ gives
\begin{equation}
k_F
=
\left(\frac{6\pi^2 n_{\mathrm{id}}}{d}\right)^{1/3}.
\end{equation}

This is exactly the formula used in the code.

\subsection{Relativistic Single-Particle Energy}

The model uses the relativistic dispersion relation
\begin{equation}
E(k)=\sqrt{(\hbar c\,k)^2+m_N^2},
\end{equation}
where $m_N$ is the nucleon mass.

\subsection{Ideal Energy Density}

The ideal energy density is
\begin{equation}
\epsilon_{\mathrm{id}}
=
\frac{d}{(2\pi)^3}\int_{|\mathbf{k}|\le k_F} E(k)\,d^3k
=
\frac{d}{2\pi^2}\int_0^{k_F} k^2 E(k)\,dk.
\end{equation}

\subsection{Ideal Pressure}

The relativistic ideal-gas pressure is
\begin{equation}
p_{\mathrm{id}}
=
\frac{d}{6\pi^2}\int_0^{k_F}\frac{(\hbar c)^2 k^4}{E(k)}\,dk.
\end{equation}

These are the quantum kinetic pieces used everywhere else in the workflow.

\subsection{Degeneracy Factors}

Two degeneracy values appear in the code:
\begin{itemize}
\item $d=4$ for symmetric nuclear matter, because proton/neutron and spin-up/spin-down
are all included together;
\item $d=2$ for a single species sector, when protons and neutrons are treated
separately.
\end{itemize}

\section{Step 2: Symmetric Clausius Matter at \texorpdfstring{$T=0$}{T=0}}

\subsection{Excluded-Volume Map}

For symmetric matter, the Clausius model introduces a free-volume factor
\begin{equation}
f(n)=1-bn.
\end{equation}

The physical density is $n$, but because only the fraction $f$ of the volume
is available, the corresponding ideal-gas density is
\begin{equation}
n_{\mathrm{id}}=\frac{n}{f(n)}=\frac{n}{1-bn}.
\end{equation}

This is exactly the hard-ball analogy: if the available volume is smaller, the
effective density inside that free volume is larger.

\subsection{Symmetric Clausius Energy Density}

The total energy density is
\begin{equation}
\epsilon(n)
=
f(n)\,\epsilon_{\mathrm{id}}\!\left(\frac{n}{f(n)}\right)
-\frac{a n^2}{1+cn}.
\end{equation}

The first term is the ideal kinetic energy living in the free volume. The
second term is the attractive mean field.

\subsection{Symmetric Clausius Pressure}

The corresponding pressure is
\begin{equation}
P(n)
=
p_{\mathrm{id}}\!\left(\frac{n}{1-bn}\right)
-\frac{a n^2}{(1+cn)^2}.
\end{equation}

If $c=0$, the attractive part reduces to the simpler van der Waals-like form
$-a n^2$ in the energy density and $-a n^2$ in the pressure derivative
structure. The Clausius denominator $1+cn$ is what makes the model more
flexible.

\section{Step 3: Saturation Constraints in Symmetric Matter}

The code calibrates symmetric matter at the empirical saturation point.

\subsection{The Empirical Inputs}

The default values are
\begin{align}
n_0 &= 0.16\ \mathrm{fm}^{-3}, \\
\left.\frac{E}{A}-m_N\right|_{n_0} &= -16\ \mathrm{MeV}.
\end{align}

The pressure at saturation must also vanish:
\begin{equation}
P(n_0)=0.
\end{equation}

\subsection{Energy per Particle}

From the energy density, the energy per baryon relative to the nucleon mass is
\begin{equation}
\frac{E}{A}-m_N
=
\frac{\epsilon(n)}{n}-m_N.
\end{equation}

\subsection{How \texorpdfstring{$a$}{a} Is Obtained at Fixed \texorpdfstring{$b,c$}{b,c}}

Take the zero-pressure condition at $n=n_0$:
\begin{equation}
0
=
p_{\mathrm{id}}(n_{\mathrm{id},0})
-\frac{a n_0^2}{(1+cn_0)^2},
\end{equation}
where
\begin{equation}
n_{\mathrm{id},0}=\frac{n_0}{1-bn_0}.
\end{equation}

Solve this algebraically for $a$:
\begin{equation}
a
=
\frac{p_{\mathrm{id}}(n_{\mathrm{id},0})(1+cn_0)^2}{n_0^2}.
\end{equation}

This step is important: at fixed $b$ and $c$, the pressure condition determines
$a$ directly.

\subsection{How \texorpdfstring{$b$}{b} Is Obtained at Fixed \texorpdfstring{$c$}{c}}

Now substitute the above $a(b,c)$ into the binding condition
\begin{equation}
\frac{\epsilon(n_0)}{n_0}-m_N=-16\ \mathrm{MeV}.
\end{equation}

Using
\begin{equation}
\epsilon(n_0)
=
(1-bn_0)\,\epsilon_{\mathrm{id}}(n_{\mathrm{id},0})
-\frac{a n_0^2}{1+cn_0},
\end{equation}
we get
\begin{equation}
\frac{(1-bn_0)\epsilon_{\mathrm{id}}(n_{\mathrm{id},0})}{n_0}
-\frac{a n_0}{1+cn_0}
-m_N
=
-16\ \mathrm{MeV}.
\end{equation}

Since $a$ is already known as a function of $b$ and $c$, this becomes a single
nonlinear equation for $b$ at fixed $c$.

The solver handles it numerically by:
\begin{enumerate}
\item scanning a range in $b$;
\item looking for a sign change in the residual;
\item applying bisection.
\end{enumerate}

\subsection{Simple Mechanical Analogy}

This is similar to solving a school-level mechanics problem where one equation
gives one constant in terms of another, and then a second condition fixes the
remaining constant.

Here:
\begin{itemize}
\item the zero-pressure condition plays the role of the first equation;
\item the binding-energy condition plays the role of the second equation.
\end{itemize}

\section{Step 4: Incompressibility and the Common Clausius Parameter \texorpdfstring{$c$}{c}}

\subsection{Definition of \texorpdfstring{$K_0$}{K0}}

The incompressibility is
\begin{equation}
K_0
=
9\left.\frac{dP}{dn}\right|_{n_0}.
\end{equation}

Physically, this is the bulk stiffness of symmetric nuclear matter. A simple
analogy is a spring constant:
\begin{itemize}
\item a large $K_0$ means the matter resists compression strongly;
\item a small $K_0$ means it is softer.
\end{itemize}

\subsection{How \texorpdfstring{$c$}{c} Is Computed}

For every trial $c$:
\begin{enumerate}
\item solve the saturation problem and find $b(c)$;
\item compute $a(c)$ from the zero-pressure condition;
\item compute
\[
K_0(c)=9\left.\frac{dP}{dn}\right|_{n_0}.
\]
\end{enumerate}

Then define the residual
\begin{equation}
F_c(c)=K_0(c)-K_0^{\mathrm{target}}.
\end{equation}

The code scans $c$ over the interval
\[
0\le c\le 4.74\ \mathrm{fm}^3,
\]
finds a sign change in $F_c(c)$, and uses bisection to obtain the final
$c$.

\subsection{What Is Fixed After This Step?}

Once $c$ is known, the symmetric problem is fully fixed. The code therefore
also knows
\begin{equation}
a_{\mathrm{avg}}=a(c),\qquad b_{\mathrm{avg}}=b(c).
\end{equation}

These are called averages because they will become the center of the later
asymmetric splitting.

\section{Step 5: Asymmetric Clausius Matter}

\subsection{Proton Fraction and Species Densities}

Define the proton fraction
\begin{equation}
y=\frac{n_p}{n},
\end{equation}
so that
\begin{equation}
n_p=yn,\qquad n_n=(1-y)n.
\end{equation}

Two special cases are central:
\begin{itemize}
\item symmetric nuclear matter: $y=0.5$;
\item pure neutron matter: $y=0$.
\end{itemize}

\subsection{Species-Dependent Excluded Volume}

The model uses
\begin{align}
x_p &= b_n n_p + b_{pn} n_n, \\
x_n &= b_{pn} n_p + b_n n_n,
\end{align}
and therefore
\begin{align}
f_p &= 1-x_p, \\
f_n &= 1-x_n.
\end{align}

The corresponding ideal densities are
\begin{align}
n_{p,\mathrm{id}} &= \frac{n_p}{f_p}, \\
n_{n,\mathrm{id}} &= \frac{n_n}{f_n}.
\end{align}

\subsection{Asymmetric Energy Density}

The total energy density is
\begin{equation}
\epsilon(n,y)
=
f_p\,\epsilon_p^{\mathrm{id}}(n_{p,\mathrm{id}})
+
f_n\,\epsilon_n^{\mathrm{id}}(n_{n,\mathrm{id}})
-\frac{a_n(n_p^2+n_n^2)+2a_{pn}n_pn_n}{1+cn}.
\end{equation}

The meaning of the coefficients is:
\begin{itemize}
\item $a_n$: same-species attractive channel;
\item $a_{pn}$: proton-neutron attractive channel;
\item $b_n$: same-species excluded-volume channel;
\item $b_{pn}$: proton-neutron excluded-volume channel;
\item $c$: common density-softening parameter in the attractive denominator.
\end{itemize}

\section{Step 6: Why the Split Is Written in Terms of Averages and Differences}

The code writes
\begin{align}
a_n &= a_{\mathrm{avg}}-\Delta a, \\
a_{pn} &= a_{\mathrm{avg}}+\Delta a, \\
b_n &= b_{\mathrm{avg}}-\Delta b, \\
b_{pn} &= b_{\mathrm{avg}}+\Delta b.
\end{align}

This is not just notation. It is chosen so that the model automatically reduces
to the symmetric fit at $y=0.5$.

\subsection{Proof for the Excluded-Volume Part}

At $y=0.5$,
\[
n_p=n_n=\frac{n}{2}.
\]
Then
\begin{align}
x_p
&=
b_n\frac{n}{2}+b_{pn}\frac{n}{2}
=
\frac{b_n+b_{pn}}{2}n
=
b_{\mathrm{avg}}n, \\
x_n
&=
b_{pn}\frac{n}{2}+b_n\frac{n}{2}
=
b_{\mathrm{avg}}n.
\end{align}

So
\begin{equation}
f_p=f_n=1-b_{\mathrm{avg}}n,
\end{equation}
which is exactly the symmetric free-volume factor.

\subsection{Proof for the Attractive Part}

At $y=0.5$,
\[
n_p=n_n=\frac{n}{2}.
\]
Then the attractive numerator becomes
\begin{align}
a_n(n_p^2+n_n^2)+2a_{pn}n_pn_n
&=
a_n\left(\frac{n^2}{4}+\frac{n^2}{4}\right)
+2a_{pn}\frac{n^2}{4} \\
&=
\frac{a_n+a_{pn}}{2}n^2 \\
&=
a_{\mathrm{avg}}n^2.
\end{align}

So the asymmetric model reduces exactly to the symmetric Clausius form.

\section{Step 7: Symmetry Energy, \texorpdfstring{$J$}{J}, and \texorpdfstring{$L$}{L}}

\subsection{Energy Per Particle}

For any density and proton fraction,
\begin{equation}
E(n,y)-m_N
=
\frac{\epsilon(n,y)}{n}-m_N.
\end{equation}

\subsection{Parabolic Symmetry-Energy Convention}

The code adopts
\begin{equation}
S(n)=E_{\mathrm{PNM}}(n)-E_{\mathrm{SNM}}(n),
\end{equation}
where
\begin{align}
E_{\mathrm{PNM}}(n) &= E(n,y=0), \\
E_{\mathrm{SNM}}(n) &= E(n,y=0.5).
\end{align}

This is called the parabolic convention because in many nuclear-matter
expansions one writes
\[
E(n,\delta)\approx E(n,0)+S(n)\delta^2,
\]
with isospin asymmetry $\delta=(n_n-n_p)/n=1-2y$. In that picture,
$S(n)$ is the coefficient of the quadratic penalty for moving away from the
symmetric state.

\subsection{The Symmetry Energy at Saturation}

The code imposes
\begin{equation}
J=S(n_0)=32.5\ \mathrm{MeV}.
\end{equation}

Physically, $J$ measures how expensive it is, in energy per particle, to turn
symmetric matter at saturation into pure neutron matter.

\subsection{The Symmetry Slope}

The slope is defined as
\begin{equation}
L
=
3n_0\left.\frac{dS}{dn}\right|_{n_0}.
\end{equation}

The code imposes
\begin{equation}
L=58.9\ \mathrm{MeV}.
\end{equation}

Physically, $L$ tells us how the asymmetry penalty changes as we compress or
expand matter around saturation density.

\section{Step 8: How Each of the Five Final Parameters Is Computed}

This is the central step of the whole workflow.

\subsection{Parameter 1: \texorpdfstring{$c$}{c}}

Choose a target incompressibility $K_0^{\mathrm{target}}$.

For each trial $c$:
\begin{enumerate}
\item solve the symmetric saturation problem;
\item compute the resulting $K_0(c)$;
\item solve
\[
F_c(c)=K_0(c)-K_0^{\mathrm{target}}=0.
\]
\end{enumerate}

The root of that equation is the final $c$.

\subsection{Intermediate Quantities: \texorpdfstring{$a_{\mathrm{avg}}$}{aavg} and \texorpdfstring{$b_{\mathrm{avg}}$}{bavg}}

Once $c$ is fixed:
\begin{itemize}
\item $b_{\mathrm{avg}}$ is the solution of the symmetric binding equation;
\item $a_{\mathrm{avg}}$ is computed from the zero-pressure formula
\[
a_{\mathrm{avg}}
=
\frac{p_{\mathrm{id}}(n_{\mathrm{id},0})(1+cn_0)^2}{n_0^2},
\qquad
n_{\mathrm{id},0}=\frac{n_0}{1-b_{\mathrm{avg}}n_0}.
\]
\end{itemize}

These are not part of the final five-parameter list, but they are the center
from which the final four asymmetric parameters are built.

\subsection{Parameter 2 and Parameter 3: \texorpdfstring{$a_n$}{an} and \texorpdfstring{$a_{pn}$}{apn}}

For a trial $\Delta b$, define
\begin{equation}
F_J(\Delta a;\Delta b)
=
J(\Delta a,\Delta b)-J_{\mathrm{target}}.
\end{equation}

The code solves
\begin{equation}
F_J(\Delta a;\Delta b)=0
\end{equation}
for $\Delta a$.

Once $\Delta a$ is known,
\begin{align}
a_n &= a_{\mathrm{avg}}-\Delta a, \\
a_{pn} &= a_{\mathrm{avg}}+\Delta a.
\end{align}

So:
\begin{itemize}
\item $a_n$ is not solved directly;
\item $a_{pn}$ is not solved directly;
\item both come from solving for the auxiliary variable $\Delta a$ through the
$J$ constraint.
\end{itemize}

\subsection{Parameter 4 and Parameter 5: \texorpdfstring{$b_n$}{bn} and \texorpdfstring{$b_{pn}$}{bpn}}

Now define
\begin{equation}
F_L(\Delta b)
=
L\big(\Delta a(\Delta b),\Delta b\big)-L_{\mathrm{target}},
\end{equation}
where $\Delta a(\Delta b)$ means that for each trial $\Delta b$, the code
first solves the $J$ constraint again.

The code then solves
\begin{equation}
F_L(\Delta b)=0
\end{equation}
for $\Delta b$.

Once $\Delta b$ is known,
\begin{align}
b_n &= b_{\mathrm{avg}}-\Delta b, \\
b_{pn} &= b_{\mathrm{avg}}+\Delta b.
\end{align}

So:
\begin{itemize}
\item $b_n$ is not solved directly;
\item $b_{pn}$ is not solved directly;
\item both come from solving for the auxiliary variable $\Delta b$ through the
$L$ constraint, while $J$ is re-enforced at each trial.
\end{itemize}

\subsection{The Full Hierarchy in One Box}

The full fitting hierarchy is:
\begin{align}
K_0 &\Longrightarrow c, \\
c &\Longrightarrow (a_{\mathrm{avg}}, b_{\mathrm{avg}}), \\
J &\Longrightarrow \Delta a, \\
L &\Longrightarrow \Delta b, \\
(\Delta a,\Delta b)
&\Longrightarrow
(a_n,a_{pn},b_n,b_{pn}).
\end{align}

That is the cleanest mathematical summary of how the five final parameters are
obtained.

\section{Step 8A: Worked Zero-Temperature Example for \texorpdfstring{$K_0=280$}{K0=280}}

It is useful to see one concrete fitted case using actual numbers from the
repository output table.

\subsection{Fitted Symmetric Core}

For the target
\[
K_0^{\mathrm{target}}=280\ \mathrm{MeV},
\]
the saved fit gives
\begin{align}
c &= 4.11410812010825\ \mathrm{fm}^3, \\
a_{\mathrm{avg}} &= 454.41269746338673\ \mathrm{MeV\ fm}^3, \\
b_{\mathrm{avg}} &= 1.9364497235562026\ \mathrm{fm}^3.
\end{align}

At saturation density,
\[
n_0=0.16\ \mathrm{fm}^{-3},
\]
the symmetric free-volume factor is
\begin{equation}
f(n_0)=1-b_{\mathrm{avg}}n_0
=
1-(1.9364497235562026)(0.16)
=
0.6901680442310076.
\end{equation}

Therefore the corresponding ideal density is
\begin{equation}
n_{\mathrm{id},0}
=
\frac{n_0}{f(n_0)}
=
\frac{0.16}{0.6901680442310076}
=
0.2318275981297767\ \mathrm{fm}^{-3}.
\end{equation}

The final output then confirms that the fitted symmetric point satisfies
\begin{align}
K_0 &= 280.00000000183434\ \mathrm{MeV}, \\
\left.\frac{E}{A}-m_N\right|_{n_0}
&=
-15.999999999746478\ \mathrm{MeV}, \\
n_{\mathrm{sat}} &= 0.16000002779816316\ \mathrm{fm}^{-3}, \\
P_{\mathrm{sat}} &= 8.648315237280713\times 10^{-7}\ \mathrm{MeV\ fm}^{-3}.
\end{align}

So the saturation and incompressibility constraints are being matched to very
high numerical accuracy.

\subsection{Asymmetric Splits}

The same fitted row gives
\begin{align}
\Delta a &= 204.076171451012\ \mathrm{MeV\ fm}^3, \\
\Delta b &= 0.5190264752342095\ \mathrm{fm}^3.
\end{align}

Using
\begin{align}
a_n &= a_{\mathrm{avg}}-\Delta a, \\
a_{pn} &= a_{\mathrm{avg}}+\Delta a, \\
b_n &= b_{\mathrm{avg}}-\Delta b, \\
b_{pn} &= b_{\mathrm{avg}}+\Delta b,
\end{align}
we obtain
\begin{align}
a_n
&=
454.41269746338673-204.076171451012 \\
&=
250.33652601237475\ \mathrm{MeV\ fm}^3, \\
a_{pn}
&=
454.41269746338673+204.076171451012 \\
&=
658.4888689143987\ \mathrm{MeV\ fm}^3,
\end{align}
and
\begin{align}
b_n
&=
1.9364497235562026-0.5190264752342095 \\
&=
1.417423248321993\ \mathrm{fm}^3, \\
b_{pn}
&=
1.9364497235562026+0.5190264752342095 \\
&=
2.455476198790412\ \mathrm{fm}^3.
\end{align}

\subsection{Consistency Checks}

The split reproduces the averages exactly:
\begin{align}
\frac{a_n+a_{pn}}{2}
&=
\frac{250.33652601237475+658.4888689143987}{2} \\
&=
454.41269746338673
=
a_{\mathrm{avg}},
\end{align}
and
\begin{align}
\frac{b_n+b_{pn}}{2}
&=
\frac{1.417423248321993+2.455476198790412}{2} \\
&=
1.9364497235562026
=
b_{\mathrm{avg}}.
\end{align}

The same row also satisfies the isovector constraints:
\begin{align}
J &= 32.499999996661245\ \mathrm{MeV}, \\
L &= 58.89999995170001\ \mathrm{MeV}.
\end{align}

So this one row is a concrete numerical realization of the full chain
\[
K_0 \Longrightarrow c \Longrightarrow
(a_{\mathrm{avg}},b_{\mathrm{avg}})
\Longrightarrow
(\Delta a,\Delta b)
\Longrightarrow
(a_n,a_{pn},b_n,b_{pn}).
\]

\section{Step 9: Numerical Methods Used by the Zero-Temperature Solver}

\subsection{Why Scan Plus Bisection?}

The code chooses a robust one-dimensional strategy:
\begin{itemize}
\item scan the parameter over a range;
\item look for a sign change in the residual;
\item apply bisection.
\end{itemize}

This is slower than a clever multidimensional Newton solve, but it is safer and
easier to interpret.

\subsection{Where Derivatives Appear}

Derivatives are used only for physical observables such as
\begin{equation}
K_0=9\left.\frac{dP}{dn}\right|_{n_0},
\qquad
L=3n_0\left.\frac{dS}{dn}\right|_{n_0},
\end{equation}
and are computed numerically with centered finite differences.

\subsection{Saturation Check}

After the symmetric point is solved, the code minimizes
\[
\frac{E}{A}-m_N
\]
in a density interval around saturation as a consistency check. This yields
\[
n_{\mathrm{sat}},\qquad
\left(\frac{E}{A}-m_N\right)_{\mathrm{sat}},\qquad
P_{\mathrm{sat}}.
\]

These are diagnostic outputs that tell us whether the fitted point really sits
where it should.

\section{Step 10: The Finite-Temperature Quantum Solver}

The second part of the repository uses the fitted five parameters to compute
the critical point at fixed proton fraction $y$.

\subsection{Finite-Temperature Ideal Fermi Gas}

At nonzero temperature, the sharp Fermi sphere is replaced by the Fermi-Dirac
occupation factor
\begin{equation}
f_{\mathrm{FD}}(E;\mu,T)=\frac{1}{e^{(E-\mu)/T}+1}.
\end{equation}

Then
\begin{equation}
n_{\mathrm{id}}(T,\mu)
=
\frac{d}{2\pi^2}\int_0^\infty k^2 f_{\mathrm{FD}}(E;\mu,T)\,dk,
\end{equation}
and
\begin{equation}
p_{\mathrm{id}}(T,\mu)
=
\frac{d}{6\pi^2}\int_0^\infty
\frac{(\hbar c)^2k^4}{E(k)}f_{\mathrm{FD}}(E;\mu,T)\,dk.
\end{equation}

This is why the finite-temperature stage is called quantum. It is using Fermi
statistics rather than a classical Boltzmann gas.

\subsection{Species Target Densities at Fixed \texorpdfstring{$y$}{y}}

For fixed proton fraction $y$:
\begin{align}
n_p &= yn, \\
n_n &= (1-y)n, \\
f_p &= 1-b_n n_p - b_{pn} n_n, \\
f_n &= 1-b_{pn} n_p - b_n n_n.
\end{align}

The ideal target densities are therefore
\begin{align}
n_{p,\mathrm{id}}^{\mathrm{target}} &= \frac{n_p}{f_p}, \\
n_{n,\mathrm{id}}^{\mathrm{target}} &= \frac{n_n}{f_n}.
\end{align}

\subsection{Inner SCF: Solving for \texorpdfstring{$\mu_p^\ast$}{mu*p} and \texorpdfstring{$\mu_n^\ast$}{mu*n}}

The solver must find effective chemical potentials $\mu_p^\ast$ and
$\mu_n^\ast$ such that
\begin{align}
n_{\mathrm{id}}(T,\mu_p^\ast)
&=
n_{p,\mathrm{id}}^{\mathrm{target}}, \\
n_{\mathrm{id}}(T,\mu_n^\ast)
&=
n_{n,\mathrm{id}}^{\mathrm{target}}.
\end{align}

The code does this with a damped Newton-like iteration.

This is similar to a standard introductory numerical-analysis problem:
\[
\text{given } F(x)=0,\ \text{update } x \text{ using the local slope of } F.
\]

\subsection{Finite-Temperature Pressure Branch}

Once $\mu_p^\ast$ and $\mu_n^\ast$ are found, the total pressure is
\begin{equation}
P(T,n;y)
=
p_p^{\mathrm{id}}(T,\mu_p^\ast)
+
p_n^{\mathrm{id}}(T,\mu_n^\ast)
-\frac{C(y)n^2}{(1+cn)^2},
\end{equation}
where
\begin{equation}
C(y)
=
a_n\big[y^2+(1-y)^2\big]
+2a_{pn}y(1-y).
\end{equation}

At $y=0.5$, one has
\begin{equation}
C(0.5)=a_{\mathrm{avg}},
\end{equation}
so the fixed-$y$ solver reduces back to the symmetric Clausius branch.

\subsection{Critical-Point Conditions}

The liquid-gas critical point is defined by the usual inflection-point
conditions on the isotherm:
\begin{align}
\left(\frac{\partial P}{\partial n}\right)_T &= 0, \\
\left(\frac{\partial^2 P}{\partial n^2}\right)_T &= 0.
\end{align}

This is the same mathematics as the standard van der Waals critical point. On
a graph of $P$ versus $n$ at fixed $T$, the critical isotherm is flat enough to
have zero slope and zero curvature at the critical point.

\subsection{Outer SCF: Solving for \texorpdfstring{$T_c$}{Tc} and \texorpdfstring{$n_c$}{nc}}

The outer solver seeks the pair $(T,n)$ that satisfies the two equations above.
It proceeds by:
\begin{enumerate}
\item scanning a coarse grid in $(T,n)$;
\item selecting a good seed;
\item updating $T$ and $n$ with damped slope-based corrections;
\item accepting only trial steps that reduce the residual score.
\end{enumerate}

So the finite-temperature code has two nested loops:
\begin{itemize}
\item an inner loop for $\mu_p^\ast$ and $\mu_n^\ast$;
\item an outer loop for $(T_c,n_c)$.
\end{itemize}

\section{Step 10A: Worked Fixed-\texorpdfstring{$y$}{y} Example for \texorpdfstring{$K_0=280$}{K0=280}}

Now take the same fitted interaction and follow it into the finite-temperature
solver.

\subsection{Symmetric Check at \texorpdfstring{$y=0.5$}{y=0.5}}

For the $K_0=280$ row, the finite-temperature solver returns at $y=0.5$
\begin{align}
T_c &= 16.4817165975974\ \mathrm{MeV}, \\
n_c &= 0.0520814359187431\ \mathrm{fm}^{-3}, \\
P_c &= 0.261959538318564\ \mathrm{MeV\ fm}^{-3}.
\end{align}

At this point,
\[
n_p=n_n=\frac{n_c}{2}.
\]
Using the fitted excluded-volume parameters,
\begin{align}
x_p
&=
b_n\frac{n_c}{2}+b_{pn}\frac{n_c}{2}
=
\frac{b_n+b_{pn}}{2}n_c \\
&=
b_{\mathrm{avg}}n_c
=
0.10085308218726015,
\end{align}
and similarly
\[
x_n=0.10085308218726015.
\]
Hence
\[
f_p=f_n=1-0.10085308218726015=0.8991469178127398.
\]

For the attractive coefficient,
\begin{align}
C(0.5)
&=
a_n\left[\left(\frac12\right)^2+\left(\frac12\right)^2\right]
+2a_{pn}\left(\frac12\right)\left(\frac12\right) \\
&=
\frac{a_n+a_{pn}}{2}
=
a_{\mathrm{avg}}
=
454.41269746338673\ \mathrm{MeV\ fm}^3.
\end{align}

So the fixed-$y$ formulas reduce exactly to the symmetric Clausius formulas,
as they should.

\subsection{Asymmetric Example at \texorpdfstring{$y=0.3$}{y=0.3}}

For the same $K_0=280$ interaction, the finite-temperature solver gives at
$y=0.3$
\begin{align}
T_c &= 14.063530458374572\ \mathrm{MeV}, \\
n_c &= 0.048943415460203074\ \mathrm{fm}^{-3}, \\
P_c &= 0.2110688843825943\ \mathrm{MeV\ fm}^{-3}, \\
\mu_p^\ast &= 934.7468901046402\ \mathrm{MeV}, \\
\mu_n^\ast &= 950.2022651752985\ \mathrm{MeV}.
\end{align}

Now compute the proton and neutron densities:
\begin{align}
n_p &= yn_c = (0.3)(0.048943415460203074)
= 0.014683024638060921\ \mathrm{fm}^{-3}, \\
n_n &= (1-y)n_c = (0.7)(0.048943415460203074)
= 0.03426039082214215\ \mathrm{fm}^{-3}.
\end{align}

Then
\begin{align}
x_p
&=
b_n n_p + b_{pn} n_n \\
&=
(1.417423248321993)(0.014683024638060921)
+(2.455476198790412)(0.03426039082214215) \\
&=
0.10493763470269968,
\end{align}
while
\begin{align}
x_n
&=
b_{pn} n_p + b_n n_n \\
&=
(2.455476198790412)(0.014683024638060921)
+(1.417423248321993)(0.03426039082214215) \\
&=
0.08461529197291352.
\end{align}

Therefore
\begin{align}
f_p &= 1-x_p = 0.8950623652973003, \\
f_n &= 1-x_n = 0.9153847080270865.
\end{align}

The corresponding ideal target densities are
\begin{align}
n_{p,\mathrm{id}}^{\mathrm{target}}
&=
\frac{n_p}{f_p}
=
\frac{0.014683024638060921}{0.8950623652973003} \\
&=
0.016404471026087514\ \mathrm{fm}^{-3},
\end{align}
and
\begin{align}
n_{n,\mathrm{id}}^{\mathrm{target}}
&=
\frac{n_n}{f_n}
=
\frac{0.03426039082214215}{0.9153847080270865} \\
&=
0.03742731391699016\ \mathrm{fm}^{-3}.
\end{align}

\subsection{Effective Attractive Coefficient at \texorpdfstring{$y=0.3$}{y=0.3}}

The fixed-$y$ attractive coefficient is
\begin{align}
C(y)
&=
a_n\big[y^2+(1-y)^2\big]+2a_{pn}y(1-y).
\end{align}

At $y=0.3$,
\begin{align}
y^2+(1-y)^2 &= 0.3^2+0.7^2 = 0.58, \\
2y(1-y) &= 2(0.3)(0.7) = 0.42.
\end{align}
Hence
\begin{align}
C(0.3)
&=
(250.33652601237475)(0.58)
+(658.4888689143987)(0.42) \\
&=
421.76051003122484\ \mathrm{MeV\ fm}^3.
\end{align}

The interaction pressure contribution at the critical density is therefore
\begin{align}
P_{\mathrm{int}}
&=
-\frac{C(0.3)n_c^2}{(1+cn_c)^2} \\
&=
-\frac{(421.76051003122484)(0.048943415460203074)^2}
{\big(1+(4.11410812010825)(0.048943415460203074)\big)^2} \\
&=
-0.7000179980681587\ \mathrm{MeV\ fm}^{-3}.
\end{align}

Since the full pressure is
\[
P_c=0.2110688843825943\ \mathrm{MeV\ fm}^{-3},
\]
the proton-plus-neutron ideal pressures together must contribute
\begin{align}
p_p^{\mathrm{id}}+p_n^{\mathrm{id}}
&=
P_c-P_{\mathrm{int}} \\
&=
0.2110688843825943-(-0.7000179980681587) \\
&=
0.911086882450753\ \mathrm{MeV\ fm}^{-3}.
\end{align}

This is a useful physical interpretation:
\begin{itemize}
\item the quantum ideal-gas sectors provide positive pressure;
\item the mean-field attraction provides a sizable negative contribution;
\item the final critical pressure is the small positive remainder after these
pieces compete.
\end{itemize}

\section{A Few Simple Physics Examples}

\subsection{Example 1: Why the Pair Term Looks Like \texorpdfstring{$n^2$}{n2}}

Suppose each particle can interact with all others in an average way. If there
are $N$ particles, the number of pairs is proportional to $N^2$. Dividing by
volume squared gives a density-squared scaling. This is why the attraction term
has the form
\[
\text{interaction energy density} \sim -a n^2.
\]

\subsection{Example 2: Why Excluded Volume Increases the Ideal Density}

Take a box of total volume $V$ containing particles whose cores occupy a volume
$V_{\mathrm{occ}}$. The free volume is
\[
V_{\mathrm{free}}=V-V_{\mathrm{occ}}.
\]
If the particle number is $N$, then the density inside the available volume is
\[
\frac{N}{V_{\mathrm{free}}}
\]
rather than
\[
\frac{N}{V}.
\]
That is exactly the logic behind
\[
n_{\mathrm{id}}=\frac{n}{1-bn}.
\]

\subsection{Example 3: Why \texorpdfstring{$J$}{J} and \texorpdfstring{$L$}{L} Need Two Different Splits}

Think of $J$ as the \emph{height} of an energy penalty at one special density
$n_0$, and think of $L$ as the \emph{slope} of that penalty there.

One number is not enough to independently control both a value and a slope.
That is why the code uses two splitting directions:
\begin{itemize}
\item $\Delta a$ mainly tunes the size of the asymmetry energy;
\item $\Delta b$ then adjusts how that asymmetry energy changes with density.
\end{itemize}

\section{Final Summary}

The workflow can be summarized as follows.

\begin{enumerate}
\item Build the relativistic ideal Fermi gas.
\item Add excluded-volume repulsion and Clausius attraction.
\item In symmetric matter, impose
\[
P(n_0)=0,\qquad
\left.\frac{E}{A}-m_N\right|_{n_0}=-16\ \mathrm{MeV}
\]
to solve $a$ and $b$ at fixed $c$.
\item Compute
\[
K_0=9\left.\frac{dP}{dn}\right|_{n_0}
\]
and use the target $K_0$ to determine $c$.
\item Rename the resulting symmetric parameters as
\[
a_{\mathrm{avg}},\qquad b_{\mathrm{avg}}.
\]
\item Introduce the asymmetric splitting
\[
a_n=a_{\mathrm{avg}}-\Delta a,\qquad
a_{pn}=a_{\mathrm{avg}}+\Delta a,
\]
\[
b_n=b_{\mathrm{avg}}-\Delta b,\qquad
b_{pn}=b_{\mathrm{avg}}+\Delta b.
\]
\item Solve $\Delta a$ from
\[
J=S(n_0)=32.5\ \mathrm{MeV}.
\]
\item Solve $\Delta b$ from
\[
L=3n_0\left.\frac{dS}{dn}\right|_{n_0}=58.9\ \mathrm{MeV},
\]
while re-enforcing the $J$ constraint at every trial.
\item The final five fitted interaction parameters are then
\[
a_n,\qquad a_{pn},\qquad b_n,\qquad b_{pn},\qquad c.
\]
\item Use those five parameters in the finite-temperature fixed-$y$ quantum
solver to determine
\[
T_c,\qquad n_c,\qquad P_c
\]
from
\[
\left(\frac{\partial P}{\partial n}\right)_T=0,
\qquad
\left(\frac{\partial^2 P}{\partial n^2}\right)_T=0.
\]
\end{enumerate}

That is the complete mathematical logic of the repository's Clausius workflow.

\end{document}
